\SourceLeaf{82}
\noindent\hspace*{\parindent}Sei $\mathfrak{P}$ die gesuchte Funktion, so gilt:

\[
\begin{aligned}\left(\frac{d\mathfrak{P}}{dP}\right) ={}& \text{konstanter Teil oder der von } q \text{ freie Teil in }\\ &\frac{q\,.\,f'q\,.\,f\dfrac{P^{\frac1n}}{Q}}{nP\left(1-\left(\dfrac{P}{Q}\right)^{\frac{\nu}{n}}\right)}.
\end{aligned}
\]



\Needspace{28\baselineskip}\begin{center}
BEMERKUNGEN ZU DEN NOTIZEN ÜBER REIHENINVERSION.
\end{center}

\noindent\hspace*{\parindent}Die erste Notiz ist einem Brief an \textsc{Olbers} vom 12.\,Dec. 1813 entnommen. Der dort mit $(n,\nu)$ bezeichnete Entwicklungskoeffizient ergibt sich aus dem Multinomialsatz explizit in der Form:

\[
(n,\nu) = \sum_\lambda \frac{(n+\mu-1)(n+\mu-2).\,.\,.(n+1)n}{\lambda'!\,\lambda''!\,\lambda'''!\,.\,.\,.\,.}\, a^{\prime\lambda'}a^{\prime\prime\lambda''}a^{\prime\prime\prime\lambda'''}.\,.\,.,
\]

\noindent wenn man dabei $a'=-\alpha$, $a''=-\beta$, $a'''=-\gamma$, .\,.\,. setzt. Die Summe erstreckt sich über alle Systeme nichtnegativer ganzer Zahlen $\lambda'$, $\lambda''$, .\,.\,., die die Gleichung $\lambda'+2\lambda''+3\lambda'''+.\,.\,.=\nu$ erfüllen, und es gilt $\mu=\lambda'+\lambda''+\lambda'''+.\,.\,.$.

Die zweite Notiz, die \textsc{Gauss} in sein Exemplar von »\textsc{Hobert} und \textsc{Ideler}, trigonometrische Tafeln« (erste Seite des Umschlags) geschrieben hat, gibt die am Schluss der ersten Notiz angedeutete Verallgemeinerung explizit an. Die fragliche Inversionsformel lässt sich mithilfe des eben hergeleiteten Ausdrucks für $(n,\nu)$ leicht bestätigen.

In der dritten Notiz, die sich auf einem einzelnen Zettel findet und das Datum 1820 Jun.\,22 trägt, wird die Verallgemeinerung der ersten Notiz für beliebige Potenzen von $x$ mit einer geringfügig abgewandelten Methode ebenfalls explizit hergeleitet. Der Schluss von den beiden Darstellungen von $\Omega$ auf die Gleichungen I, II, .\,.\,. beruht auf einer Homogenitätsbetrachtung.

Die vierte Entwicklung stammt vermutlich aus dem Jahr 1816 und findet sich in einem Handbuch. Das Symbol $[Fx]_{\text{Coeff.}\,x^n}$ bezeichnet hier den Koeffizienten von $x^n$ in der Entwicklung von $F(x)$ nach Potenzen von $x$. \hfill\textsc{Fricke.}

